\documentclass[UTF8]{article}
                \usepackage[UTF8]{ctex}
                \usepackage{amssymb}
                \usepackage{geometry}
                \usepackage{array}
                \usepackage{multirow}
                \usepackage{setspace}
                \usepackage{lastpage}
                \usepackage{fancyhdr}
                \usepackage{graphicx}
                \usepackage{ragged2e}
                \usepackage{xcolor} 
                \usepackage{datetime} 
                \usepackage{advdate} 

                % ===== 统一命名的占位符 =====
                \newcommand{\studentid}{ 211300032}  % 学号
\newcommand{\studentname}{ 何翊豪}  % 姓名
\newcommand{\username}{ 何翊豪 }  % 用户名
\newcommand{\papertitle}{一种基于表征学习的高效持续强化学习方法}  % 论文标题

\newcommand{\formatScoreOneA}{\checkmark}\newcommand{\formatScoreOneB}{}\newcommand{\formatScoreOneC}{}\newcommand{\formatScoreOneD}{}
\newcommand{\formatScoreTwoA}{}\newcommand{\formatScoreTwoB}{\checkmark}\newcommand{\formatScoreTwoC}{}\newcommand{\formatScoreTwoD}{}
\newcommand{\formatScoreThreeA}{\checkmark}\newcommand{\formatScoreThreeB}{}\newcommand{\formatScoreThreeC}{}\newcommand{\formatScoreThreeD}{}
\newcommand{\formatScoreFourA}{\checkmark}\newcommand{\formatScoreFourB}{}\newcommand{\formatScoreFourC}{}\newcommand{\formatScoreFourD}{}
\newcommand{\formatScoreFiveA}{}\newcommand{\formatScoreFiveB}{}\newcommand{\formatScoreFiveC}{\checkmark}\newcommand{\formatScoreFiveD}{}
\newcommand{\formatScoreSixA}{\checkmark}\newcommand{\formatScoreSixB}{}\newcommand{\formatScoreSixC}{}\newcommand{\formatScoreSixD}{}
\newcommand{\topicScoreOneA}{\checkmark}\newcommand{\topicScoreOneB}{}\newcommand{\topicScoreOneC}{}\newcommand{\topicScoreOneD}{}
\newcommand{\topicScoreTwoA}{\checkmark}\newcommand{\topicScoreTwoB}{}\newcommand{\topicScoreTwoC}{}\newcommand{\topicScoreTwoD}{}
\newcommand{\topicScoreThreeA}{\checkmark}\newcommand{\topicScoreThreeB}{}\newcommand{\topicScoreThreeC}{}\newcommand{\topicScoreThreeD}{}
\newcommand{\levelScoreOneA}{}\newcommand{\levelScoreOneB}{\checkmark}\newcommand{\levelScoreOneC}{}\newcommand{\levelScoreOneD}{}
\newcommand{\levelScoreTwoA}{}\newcommand{\levelScoreTwoB}{\checkmark}\newcommand{\levelScoreTwoC}{}\newcommand{\levelScoreTwoD}{}
\newcommand{\levelScoreThreeA}{}\newcommand{\levelScoreThreeB}{\checkmark}\newcommand{\levelScoreThreeC}{}\newcommand{\levelScoreThreeD}{}
\newcommand{\levelScoreFourA}{}\newcommand{\levelScoreFourB}{\checkmark}\newcommand{\levelScoreFourC}{}\newcommand{\levelScoreFourD}{}
\newcommand{\levelScoreFiveA}{}\newcommand{\levelScoreFiveB}{\checkmark}\newcommand{\levelScoreFiveC}{}\newcommand{\levelScoreFiveD}{}
\newcommand{\levelScoreSixA}{\checkmark}\newcommand{\levelScoreSixB}{}\newcommand{\levelScoreSixC}{}\newcommand{\levelScoreSixD}{}
\newcommand{\qualityScoreOneA}{}\newcommand{\qualityScoreOneB}{\checkmark}\newcommand{\qualityScoreOneC}{}\newcommand{\qualityScoreOneD}{}
\newcommand{\qualityScoreTwoA}{}\newcommand{\qualityScoreTwoB}{\checkmark}\newcommand{\qualityScoreTwoC}{}\newcommand{\qualityScoreTwoD}{}
\newcommand{\qualityScoreThreeA}{}\newcommand{\qualityScoreThreeB}{\checkmark}\newcommand{\qualityScoreThreeC}{}\newcommand{\qualityScoreThreeD}{}
\newcommand{\totalScore}{86}  % 总评分
\newcommand{\suggestions}{ \par[关键词]关键词之间应使用"；"分开\par \par \par \par[关键词]关键词之间应使用"；"分开\par[第一章 绪论] 1.1节: 增加对持续强化学习主流方法的系统分类和代表性算法介绍。\par[第二章 相关工作] 2.2.4节: 具体化ESCP和RESeL的局限性描述，突出改进空间。\par[第三章 本文方法] 全章: 补充关键超参数设置和网络结构细节，增强方法可复现性。\par[第四章 实验] 4.1节: 补充性能提升的统计显著性检验，增强结果可信度。\par[第四章 实验] 4.2节: 明确MuJoCo实验的学术动机，避免使用'出于好奇'等表述。}  % 修改建议

% 日期
                \newcommand{\autoyear}{\the\year} 
                \newcommand{\automonth}{\ifcase\month\or 01\or 02\or 03\or 04\or 05\or 06\or 07\or 08\or 09\or 10\or 11\or 12\fi} 
                \newcommand{\autoday}{\ifnum\day<10 0\fi\number\day} 
                % =========================

                % 设置水印样式
                \definecolor{watermark}{gray}{0.75} 
                \newcommand{\watermarktext}{
                    \raisebox{1.5\baselineskip}{ 
                        \parbox{\textwidth}{
                            \centering
                            \textcolor{watermark}{%
                                \fontsize{12}{12}\selectfont%
                                本报告由睿文智评AI预审评估系统通过大语言模型生成，需人工复核后使用，不能直接作为最终评价\\%
                                生成时间为 \today\ \currenttime，用户名为 \username %
                            }%
                        }%
                    }%
                }

                % 配置页脚
                \pagestyle{fancy}
                \fancyhf{} 
                \renewcommand{\headrulewidth}{0pt} 
                \fancyfoot[C]{\watermarktext} 
                \fancyfoot[L]{}  
                \fancyfoot[R]{} 

                \newcolumntype{C}[1]{>{\centering\arraybackslash}p{#1}}
                \geometry{a4paper, left=0.5cm, right=0.5cm, top=1.5cm, bottom=1.5cm}

                \begin{document}

                \setlength{\parindent}{0pt}
                \setlength{\arrayrulewidth}{0.3pt} 
                \renewcommand{\arraystretch}{1.2}

                % 标题
                \begin{center}
                    \Large\textbf{人工智能学院本科毕设论文院内预审表}
                \end{center}

                \vspace*{1cm}

                \centering
                \setlength{\tabcolsep}{3pt} 
                \renewcommand{\arraystretch}{1.5}
                \begin{tabular}{|@{}>{}p{2.23cm}@{}|p{5.6cm}|p{2.72cm}|p{2.96cm}|@{}*{4}{>{\centering\arraybackslash}p{0.89cm}@{}|}}
                    \hline
                    学号 & \studentid & 姓名 & \multicolumn{5}{c|}{\studentname} \\
                    \hline
                    论文题目 & \multicolumn{7}{c|}{\papertitle} \\
                    \hline
                    \multicolumn{8}{|c|}{请参照评分标准，对论文打分}  \\
                    \hline
                    评议项目 & \multicolumn{3}{|c|}{评分标准} & 
                    优
                    秀 & 
                    良
                    好 & 
                    一
                    般 & 
                    较
                    差 \\
                    \hline
                    \multirow{6}{*}{论文格式} 
                    & \multicolumn{3}{p{11.28cm}|}{顺序：封面-承诺书-摘要（中、英文）-目录-正文-参考文献-致谢-附录完备，承诺书签名} & 
                    \formatScoreOneA & \formatScoreOneB & \formatScoreOneC & \formatScoreOneD \\
                    \cline{2-8}
                    & \multicolumn{3}{p{11.28cm}|}{摘要内容以300—600字为宜，摘要与关键词应在同一页；

                    关键词一般 3—5 个，摘要关键词用"；"分开；

                    英文摘要和关键词内容与中文摘要和关键词相一致，其中英文摘要以约 300 个实词为宜} & 
                    \formatScoreTwoA & \formatScoreTwoB & \formatScoreTwoC & \formatScoreTwoD \\
                    \cline{2-8}
                    & \multicolumn{3}{p{11.28cm}|}{论文目录要求标题层次清晰，目录中的标题要与正文中的标题一致} & 
                    \formatScoreThreeA & \formatScoreThreeB & \formatScoreThreeC & \formatScoreThreeD \\
                    \cline{2-8}
                    & \multicolumn{3}{p{11.28cm}|}{主体部分一般从引言（绪论）开始，以结论或讨论结束，其中引言（绪论）应独立成章；

                    主体部分每一章应另起页，一般不少于 15000 字或相当信息量（包括图表）；

                    图表、公式注意格式和编号，表的题目在表的上方，图的题目在图的下方，格式"表/图 1-3"，公式（1-3），表和图按各自顺序分开编号} & 
                    \formatScoreFourA & \formatScoreFourB & \formatScoreFourC & \formatScoreFourD \\
                    \cline{2-8}
                    & \multicolumn{3}{p{11.28cm}|}{参考文献表应置于正文后并另起页；

                    所有被引用文献均要列入参考文献表中；

                    引文采用著作-出版年制标注时，参考文献表应按著者字顺和出版年排序；

                    文中上标，页面脚注，参考文献格式符合学术标准，左顶格，[1]格式标记} & 
                    \formatScoreFiveA & \formatScoreFiveB & \formatScoreFiveC & \formatScoreFiveD \\
                    \cline{2-8}
                    & \multicolumn{3}{p{11.28cm}|}{谢辞应以简短的文字表示自己的谢意，内容限一页。} & 
                    \formatScoreSixA & \formatScoreSixB & \formatScoreSixC & \formatScoreSixD \\
                    \hline

                    \multirow{3}{*}{论文选题}
                    & \multicolumn{3}{p{11.28cm}|}{符合本学科专业培养目标，达到科学研究和实践能力培养的目的} & 
                    \topicScoreOneA & \topicScoreOneB & \topicScoreOneC & \topicScoreOneD \\
                    \cline{2-8}
                    & \multicolumn{3}{p{11.28cm}|}{满足专业培养方案中对素质、能力和知识结构的要求，工作量适当} & 
                    \topicScoreTwoA & \topicScoreTwoB & \topicScoreTwoC & \topicScoreTwoD \\
                    \cline{2-8}
                    & \multicolumn{3}{p{11.28cm}|}{选题符合本学科专业的发展，具有一定的科技或应用的参考价值} & 
                    \topicScoreThreeA & \topicScoreThreeB & \topicScoreThreeC & \topicScoreThreeD \\
                    \hline

                    \multirow{2}{*}{论文水平}
                    & \multicolumn{3}{p{11.28cm}|}{基本掌握检索中外文献资料的方法，对资料进行初步分析、综合、归纳等整理，并能适当应用。} & 
                    \levelScoreOneA & \levelScoreOneB & \levelScoreOneC & \levelScoreOneD \\
                    \cline{2-8}
                    & \multicolumn{3}{p{11.28cm}|}{能够综合应用所学知识，对课题所研究问题进行分析，研究目标明确，内容具体，且具有一定的深度。} & 
                    \levelScoreTwoA & \levelScoreTwoB & \levelScoreTwoC & \levelScoreTwoD \\
                    \hline
                \end{tabular}


                \centering
                \setlength{\tabcolsep}{3pt} 
                \renewcommand{\arraystretch}{1.5}
                \begin{tabular}{|@{}>{}p{2.23cm}@{}|p{5.6cm}|p{2.72cm}|p{2.96cm}|@{}*{4}{>{\centering\arraybackslash}p{0.89cm}@{}|}}
                    \hline
                    \multirow{4}{*}{}
                    & \multicolumn{3}{p{11.28cm}|}{较熟练运用本专业设计或研究的方法、手段和工具开展课题的设计与研究工作} & 
                    \levelScoreThreeA & \levelScoreThreeB & \levelScoreThreeC & \levelScoreThreeD \\
                    \cline{2-8}
                    & \multicolumn{3}{p{11.28cm}|}{已基本掌握了专业技能和研究方法，有一定的实践能力和水平} & 
                    \levelScoreFourA & \levelScoreFourB & \levelScoreFourC & \levelScoreFourD \\
                    \cline{2-8}
                    & \multicolumn{3}{p{11.28cm}|}{(1)熟练使用软件完成论文的录入、排版，质量较高 

                    (2)能选用专业软件或相应软件进行编程或建模、分析等工作；编程或软件使用水平较高

                    (3)外文摘要能概括论文的主要内容，用词较准确，语法较规范；能查阅并引用本专业外文文献
                    } & 
                    \levelScoreFiveA & \levelScoreFiveB & \levelScoreFiveC & \levelScoreFiveD \\
                    \cline{2-8}
                    & \multicolumn{3}{p{11.28cm}|}{(1)论文：基于选题的研究现状，进行科学的分析与综合，提出新问题，探索解决问题的方法、手段有一定的特色或新意；或有新见解 
                    (2)设计：将专业知识、技能应用于实际问题的解决，方法或思路有一定的特色或创新} & 
                    \levelScoreSixA & \levelScoreSixB & \levelScoreSixC & \levelScoreSixD \\
                    \hline

                    \multirow{3}{*}{论文质量}
                    & \multicolumn{3}{p{11.28cm}|}{概念清楚，内容正确，数据可靠，论据较充分，论证较严密，结论基本正确} & 
                    \qualityScoreOneA & \qualityScoreOneB & \qualityScoreOneC & \qualityScoreOneD \\
                    \cline{2-8}
                    & \multicolumn{3}{p{11.28cm}|}{能够完整地反映实际完成的工作，结构较严谨，语言通顺} & 
                    \qualityScoreTwoA & \qualityScoreTwoB & \qualityScoreTwoC & \qualityScoreTwoD \\
                    \cline{2-8}
                    & \multicolumn{3}{p{11.28cm}|}{(1)论文：有一定的学术价值 

                    (2)设计：有实物作品、实际运行的系统或具有一定复杂度的原型系统} & 
                    \qualityScoreThreeA & \qualityScoreThreeB & \qualityScoreThreeC & \qualityScoreThreeD \\
                    \hline

                    \hline
                    \multicolumn{4}{|p{13.51cm}|}{总体评价（给出百分制总评成绩：100-90为优秀；89-75为良好；74-60为合格；60分以下为不合格）} & \multicolumn{4}{c|}{\totalScore} \\
                    \hline
                    \multicolumn{8}{|p{17.05cm}|}{
                        \begin{minipage}[t][11.24cm][t]{\linewidth}
                            论文修改建议: 
                            \par
                            \vspace*{0pt}


                            \suggestions
                            \vfill
                            ~
                        \end{minipage}
                    } \\
                    \hline
                    \multicolumn{8}{|r|}{
                        \autoyear \hspace{0.5em} 年\hspace{1em}
                        \automonth \hspace{0.5em} 月\hspace{1em}
                        \autoday \hspace{0.5em} 日\hspace{1em}
                    } \\
                    \hline
                \end{tabular}
                \end{document}